\chapter{Complexity — Logspace Uniform Scan}
%%% generated by the blueprint pipeline from TCSlib.Complexity.CircuitComplexity.LogspaceUniformScan
\section{Overview}

Declarations of the Lean module \texttt{TCSlib.Complexity.CircuitComplexity.LogspaceUniformScan}.

\section{Declarations}

\begin{definition}[Adjacency Scan Tasks]
\lean{BoolCircuit.AdjScan.Task}
\label{BoolCircuit.AdjScan.Task}
\leanok
\uses{BoolCircuit.GateKind}
A scan task is one of three kinds: determine the circuit size, determine whether a gate
has a specified optional type (where the option may be absent), or determine whether a
specified vertex is an input to a gate.
\end{definition}

\begin{definition}[Adjacency Scan Labels]
\lean{BoolCircuit.AdjScan.Lbl}
\label{BoolCircuit.AdjScan.Lbl}
\leanok
The scan uses labels for its start, size-counting, gate-reading, hit-testing,
type-decoding, and accepting or rejecting states. Some labels carry one Boolean flag,
and the labels $kB$ and $kC$ carry one Boolean each while $kD$ carries two Booleans.
\end{definition}

\begin{definition}[Gate Kind Decoder]
\lean{BoolCircuit.AdjScan.decodeKind}
\label{BoolCircuit.AdjScan.decodeKind}
\leanok
\uses{BoolCircuit.GateKind}
The decoder maps two Boolean bits to a gate kind: $(0,0)$ gives AND, $(0,1)$ gives OR,
and either pair whose first bit is $1$ gives NOT.
\end{definition}

\begin{definition}[Unary-Prefix Read Instruction]
\lean{BoolCircuit.AdjScan.rd}
\label{BoolCircuit.AdjScan.rd}
\leanok
\uses{BoolCircuit.AdjScan.Lbl, Complexity.LogProg.Ins}
For any two destination labels, the instruction calls decider $0$ on the unary-prefix
mode of register $0$ and branches to the first or second label according to its Boolean
answer.
\end{definition}

\begin{definition}[Task-Dependent Hit Test]
\lean{BoolCircuit.AdjScan.hit}
\label{BoolCircuit.AdjScan.hit}
\leanok
\uses{BoolCircuit.AdjScan.Lbl, BoolCircuit.AdjScan.Task, Complexity.LogProg.Ins}
For any task and two destination labels, the hit-test instruction branches according to
whether register $1$'s binary representation equals the second inner payload when the
task is an edge query; for every other task, it branches according to whether that
representation equals the input's binary payload.
\end{definition}

\begin{definition}[Default Hit Answer]
\lean{BoolCircuit.AdjScan.nAns}
\label{BoolCircuit.AdjScan.nAns}
\leanok
\uses{BoolCircuit.AdjScan.Task}
The default hit answer is true for a size task, true for a type task precisely when its
requested type is absent, and false for an edge task.
\end{definition}

\begin{definition}[Gate Hit Answer]
\lean{BoolCircuit.AdjScan.hitAns}
\label{BoolCircuit.AdjScan.hitAns}
\leanok
\uses{BoolCircuit.AdjScan.Task, BoolCircuit.DAGGate}
Given a vertex index $i$ and a gate $g$, the hit answer is true for a size task, true
for a type task precisely when the requested optional type equals the gate's kind, and
true for an edge task precisely when $i$ occurs among the gate's argument indices.
\end{definition}

\begin{definition}[Adjacency Scan Program]
\lean{BoolCircuit.AdjScan.scan}
\label{BoolCircuit.AdjScan.scan}
\leanok
\uses{BoolCircuit.AdjScan.Lbl, BoolCircuit.AdjScan.Task, BoolCircuit.AdjScan.decodeKind, BoolCircuit.AdjScan.hit, BoolCircuit.AdjScan.nAns, BoolCircuit.AdjScan.rd, Complexity.LogProg.ARM}
For each task, the scan program is a register-machine program with three registers, one
called decider, and the listed scan labels. It first validates a
unary-prefix-and-payload input for size or type queries, or a unary-prefix input with
two canonical binary payloads for edge queries; invalid inputs are rejected. It then
counts the unary prefix in register $0$, incrementing register $1$ once per prefix
symbol, and tests the corresponding payload against that count: the default hit answer
determines the result if it matches, and otherwise the program advances the count and
reads the next gate record. A failed read rejects; a matching record accepts immediately
for size queries, checks the gate's two encoded kind bits for type queries, and for edge
queries compares the first payload with the scanned argument values, accepting if one
matches and continuing otherwise. The type check returns true exactly when the requested
optional kind equals the kind decoded from those two bits; the scan's remaining labels
implement the reads, counter updates, and these accept or reject decisions.
\end{definition}

\begin{definition}[Three-coordinate natural number vector]
\lean{BoolCircuit.AdjScan.rg}
\label{BoolCircuit.AdjScan.rg}
\leanok
For natural numbers $p$, $c$, and $a$, define a function on the three-element index set
whose values at indices $0$, $1$, and $2$ are respectively $p$, $c$, and $a$.
\end{definition}

\begin{lemma}[First Register Value]
\lean{BoolCircuit.AdjScan.rg_zero}
\label{BoolCircuit.AdjScan.rg_zero}
\leanok
\uses{BoolCircuit.AdjScan.rg}
\difficulty{1}
For natural numbers $p,c,a$, the function from the three-element index set to the
natural numbers that assigns values $p,c,a$ at indices $0,1,2$, respectively, takes the
value $p$ at index $0$.
\end{lemma}

\begin{lemma}[Middle register value]
\lean{BoolCircuit.AdjScan.rg_one}
\label{BoolCircuit.AdjScan.rg_one}
\leanok
\uses{BoolCircuit.AdjScan.rg}
\difficulty{1}
For any natural numbers $p,c,a$, the middle register of the assignment $(p,c,a)$ has
value $c$.
\end{lemma}

\begin{lemma}[Third Register Value]
\lean{BoolCircuit.AdjScan.rg_two}
\label{BoolCircuit.AdjScan.rg_two}
\leanok
\uses{BoolCircuit.AdjScan.rg}
\difficulty{1}
For natural numbers $p,c,a$, the function from the three-element index set to the
natural numbers that assigns values $p,c,a$ at indices $0,1,2$, respectively, takes the
value $a$ at index $2$.
\end{lemma}

\begin{lemma}[Cases for three indices]
\lean{BoolCircuit.AdjScan.fin3_cases}
\label{BoolCircuit.AdjScan.fin3_cases}
\leanok
\difficulty{4}
For any predicate $P$ on the three-element set $\{0,1,2\}$, if $P$ holds at $0$, $1$,
and $2$, then it holds at every element of that set.
\end{lemma}

\begin{lemma}[Updating the first register]
\lean{BoolCircuit.AdjScan.upd_zero}
\label{BoolCircuit.AdjScan.upd_zero}
\leanok
\uses{BoolCircuit.AdjScan.fin3_cases, BoolCircuit.AdjScan.rg}
\difficulty{4}
For natural numbers $p,c,a,q$, let $r:\{0,1,2\}\to\mathbb{N}$ be given by $r(0)=p$,
$r(1)=c$, and $r(2)=a$. Replacing $r(0)$ by $q$ gives the function whose values at
$0,1,2$ are respectively $q,c,a$.
\end{lemma}

\begin{lemma}[Updating the middle coordinate]
\lean{BoolCircuit.AdjScan.upd_one}
\label{BoolCircuit.AdjScan.upd_one}
\leanok
\uses{BoolCircuit.AdjScan.fin3_cases, BoolCircuit.AdjScan.rg}
\difficulty{4}
For natural numbers $p,c,a,q$, let $r:\{0,1,2\}\to\mathbb{N}$ be given by $r(0)=p$,
$r(1)=c$, and $r(2)=a$. Replacing $r(1)$ by $q$ gives the function whose values at
$0,1,2$ are respectively $p,q,a$.
\end{lemma}

\begin{lemma}[Updating the Third Register]
\lean{BoolCircuit.AdjScan.upd_two}
\label{BoolCircuit.AdjScan.upd_two}
\leanok
\uses{BoolCircuit.AdjScan.fin3_cases, BoolCircuit.AdjScan.rg}
\difficulty{4}
For natural numbers $p,c,a,q$, let $r:	ext{Fin}(3)\to\mathbb{N}$ be given by $r(0)=p$,
$r(1)=c$, and $r(2)=a$. Replacing $r(2)$ by $q$ yields the function whose values at
$0,1,2$ are respectively $p,c,q$.
\end{lemma}

\begin{lemma}[Zero-register initialization identity]
\lean{BoolCircuit.AdjScan.rg_init}
\label{BoolCircuit.AdjScan.rg_init}
\leanok
\uses{BoolCircuit.AdjScan.fin3_cases, BoolCircuit.AdjScan.rg}
\difficulty{4}
The constant function on the three registers that assigns zero to each register is the
register assignment with all three values equal to zero.
\end{lemma}

\begin{definition}[Bounded configuration counters]
\lean{BoolCircuit.AdjScan.Gd}
\label{BoolCircuit.AdjScan.Gd}
\leanok
\uses{BoolCircuit.AdjScan.Lbl, Complexity.LogProg.AConf}
A configuration is bounded by $B$ if each of its three natural-number counters is at
most $B$.
\end{definition}

\begin{lemma}[Register bound for a register triple]
\lean{BoolCircuit.AdjScan.gd_rg}
\label{BoolCircuit.AdjScan.gd_rg}
\leanok
\uses{BoolCircuit.AdjScan.Gd, BoolCircuit.AdjScan.Lbl, BoolCircuit.AdjScan.fin3_cases, BoolCircuit.AdjScan.rg}
\difficulty{4}
Let $B,p,c,a$ be natural numbers. For any control state and any answer status, if $p\le
B$, $c\le B$, and $a\le B$, then the configuration with these three register values has
every register at most $B$.
\end{lemma}

\begin{definition}[Encoded input word]
\lean{BoolCircuit.AdjScan.X}
\label{BoolCircuit.AdjScan.X}
\leanok
\uses{Turing.pairEncode}
For a natural number $n$ and a Boolean word $w$, define the encoded word to be the pair
encoding of the length-$n$ word consisting entirely of true bits together with $w$.
\end{definition}

\begin{definition}[Task target encoding]
\lean{BoolCircuit.AdjScan.Target}
\label{BoolCircuit.AdjScan.Target}
\leanok
\uses{BoolCircuit.AdjScan.Task, Complexity.LogProg.pairWords, Complexity.LogProg.plainWord}
For a task $	au$, a Boolean word $x$, and natural numbers $i$ and $t$, define the target
condition as follows: when $	au$ is the edge task, $x$ encodes the pair of binary
expansions of $i$ and $t$; for either other task, $x$ encodes the binary expansion of
$t$ as a plain word.
\end{definition}

\begin{lemma}[One-step oracle read]
\lean{BoolCircuit.AdjScan.astep_rd}
\label{BoolCircuit.AdjScan.astep_rd}
\leanok
\uses{BoolCircuit.AdjScan.Lbl, BoolCircuit.AdjScan.X, BoolCircuit.AdjScan.rd, BoolCircuit.AdjScan.rg, BoolCircuit.AdjScan.scan, Complexity.LogProg.astep, Complexity.LogProg.vword_unary₁}
\difficulty{3}
Let $n$ be a natural number, let $w,E$ be bit strings, let $o$ be an oracle, and let
$\tau$ be a task. Suppose that on every encoding of a natural number $p$, the oracle
returns the $p$th bit of $E$, defaulting to false beyond its end. If the scan
instruction at label $l$ is a read instruction with successor labels $l_1,l_0$, then one
machine step from label $l$ with registers $(p,c,a)$ leaves the registers unchanged and
moves to $l_1$ if the $p$th bit of $E$ (default false) is true, and to $l_0$ otherwise.
\end{lemma}

\begin{lemma}[One-step gate-hit transition]
\lean{BoolCircuit.AdjScan.astep_hit}
\label{BoolCircuit.AdjScan.astep_hit}
\leanok
\uses{BoolCircuit.AdjScan.Lbl, BoolCircuit.AdjScan.Target, BoolCircuit.AdjScan.X, BoolCircuit.AdjScan.hit, BoolCircuit.AdjScan.rg, BoolCircuit.AdjScan.scan, Complexity.LogProg.astep, Complexity.LogProg.bits_injective}
\difficulty{4}
Let $n$ be a natural number, let $w$ be a bit string, let $o$ be an oracle, and let
$\tau$ be a scan task. Suppose the input formed from $n$ and $w$ has target vertex
$tv$ (and source $iv$ for an edge task). If the scan instruction at label $l$ is the
task's hit instruction with successor labels $l_1,l_0$, then one machine step from label
$l$ with registers $(p,c,a)$ leaves the registers unchanged and moves to $l_1$ when
$c=tv$, and to $l_0$ otherwise.
\end{lemma}

\begin{lemma}[One-step source comparison]
\lean{BoolCircuit.AdjScan.astep_jeqFst}
\label{BoolCircuit.AdjScan.astep_jeqFst}
\leanok
\uses{BoolCircuit.AdjScan.Target, BoolCircuit.AdjScan.X, BoolCircuit.AdjScan.rg, BoolCircuit.AdjScan.scan, Complexity.LogProg.astep, Complexity.LogProg.bits_injective}
\difficulty{3}
Let $n$ be a natural number, let $w$ be a bit string, and suppose the input formed from
$n$ and $w$ has edge target $(iv,tv)$. For any oracle $o$ and register values $p,c,a$,
one machine step of the edge scan from the source-comparison state with registers
$(p,c,a)$ leaves the registers unchanged and moves to the accepting state when $a=iv$,
and to the continuation state otherwise.
\end{lemma}

\begin{lemma}[Entry after a list drop]
\lean{BoolCircuit.AdjScan.getD_of_drop}
\label{BoolCircuit.AdjScan.getD_of_drop}
\leanok
\difficulty{3}
Let $E$ and $r$ be finite bit strings, let $p$ be a natural number, and let $b$ be a
bit. If removing the first $p$ bits of $E$ gives $b$ followed by $r$, then the bit of
$E$ at position $p$, with out-of-range positions read as $0$, is $b$.
\end{lemma}

\begin{lemma}[Length bound for a nonempty list prefix]
\lean{BoolCircuit.AdjScan.le_of_drop}
\label{BoolCircuit.AdjScan.le_of_drop}
\leanok
\difficulty{5}
Let $E,s,$ and $r$ be bit lists and let $p$ be a natural number. If dropping the first
$p$ entries of $E$ gives the concatenation of $s$ and $r$, and $s$ is nonempty, then
$p+|s|\le |E|$.
\end{lemma}

\begin{lemma}[Successive list drops]
\lean{BoolCircuit.AdjScan.drop_succ_of_drop}
\label{BoolCircuit.AdjScan.drop_succ_of_drop}
\leanok
\difficulty{1}
Let $E$ and $r$ be finite bit strings, let $p$ be a natural number, and let $b$ be a
bit. If removing the first $p$ bits of $E$ gives the string beginning with $b$ and
followed by $r$, then removing the first $p+1$ bits of $E$ gives $r$.
\end{lemma}

\begin{lemma}[Successive list drops]
\lean{BoolCircuit.AdjScan.drop_add_of_drop}
\label{BoolCircuit.AdjScan.drop_add_of_drop}
\leanok
\difficulty{1}
Let $E$, $s$, and $r$ be finite bit strings, and let $p$ be a natural number. If
removing the first $p$ bits of $E$ leaves the concatenation of $s$ and $r$, then
removing the first $p+|s|$ bits of $E$ leaves $r$.
\end{lemma}

\begin{lemma}[Bit in a unary run]
\lean{BoolCircuit.AdjScan.bit_rep}
\label{BoolCircuit.AdjScan.bit_rep}
\leanok
\difficulty{4}
Suppose that after position $p$, the bit string $E$ consists of $k$ true bits, then a
false bit, then a suffix. For every $j\le k$, the bit of $E$ at position $p+j$,
defaulting to false if out of range, is true exactly when $j<k$.
\end{lemma}

\begin{lemma}[Unary Scan Loop]
\lean{BoolCircuit.AdjScan.uloop}
\label{BoolCircuit.AdjScan.uloop}
\leanok
\uses{BoolCircuit.AdjScan.Gd, BoolCircuit.AdjScan.X, BoolCircuit.AdjScan.astep_rd, BoolCircuit.AdjScan.bit_rep, BoolCircuit.AdjScan.gd_rg, BoolCircuit.AdjScan.le_of_drop, BoolCircuit.AdjScan.rg, BoolCircuit.AdjScan.rg_two, BoolCircuit.AdjScan.rg_zero, BoolCircuit.AdjScan.scan, BoolCircuit.AdjScan.upd_two, BoolCircuit.AdjScan.upd_zero, Complexity.LogProg.AReach, Complexity.LogProg.AReach.refl, Complexity.LogProg.AReach.step, Complexity.LogProg.astep}
\difficulty{6}
Let $E$ be a finite bit string of length at most $B$, and let $o$ return the bit of $E$
at each natural-number position, defaulting to false beyond its end. Suppose the suffix
of $E$ beginning at position $p$ consists of exactly $k$ true bits followed by a false
bit and a remaining suffix. For either value of the Boolean flag $ex$, the scan machine,
on input $X(n,w)$ and starting in state $uA(ex)$ with registers $(p,c,0)$ and no result,
can reach state $uE(ex)$ with registers $(p+k,c,k)$ and no result, with every register
bounded by $B$ at all intermediate configurations.
\end{lemma}

\begin{lemma}[Length of an encoded nonempty list]
\lean{BoolCircuit.AdjScan.length_encodeList_cons}
\label{BoolCircuit.AdjScan.length_encodeList_cons}
\leanok
\uses{BoolCircuit.encodeList}
\difficulty{3}
For a function $f$ from a type $\alpha$ to bit strings, an element $a\in\alpha$, and a
list $as$ of elements of $\alpha$, the encoding of $a::as$ has length
$1+|f(a)|+|\operatorname{encode}(as)|$, where the list encoding uses $f$.
\end{lemma}

\begin{lemma}[Skipping an encoded argument list]
\lean{BoolCircuit.AdjScan.argsSkip}
\label{BoolCircuit.AdjScan.argsSkip}
\leanok
\uses{BoolCircuit.AdjScan.Gd, BoolCircuit.AdjScan.X, BoolCircuit.AdjScan.astep_rd, BoolCircuit.AdjScan.drop_add_of_drop, BoolCircuit.AdjScan.drop_succ_of_drop, BoolCircuit.AdjScan.gd_rg, BoolCircuit.AdjScan.getD_of_drop, BoolCircuit.AdjScan.le_of_drop, BoolCircuit.AdjScan.length_encodeList_cons, BoolCircuit.AdjScan.rg, BoolCircuit.AdjScan.rg_zero, BoolCircuit.AdjScan.scan, BoolCircuit.AdjScan.uloop, BoolCircuit.AdjScan.upd_two, BoolCircuit.AdjScan.upd_zero, BoolCircuit.encodeList, BoolCircuit.encodeNat, Complexity.LogProg.AReach, Complexity.LogProg.AReach.refl, Complexity.LogProg.AReach.step, Complexity.LogProg.AReach.trans, Complexity.LogProg.astep}
\difficulty{6}
Let $n$ be a natural number, let $w,E$ be bit strings, and let $o$ be an oracle on bit
strings. Suppose that on every encoding of a natural number $p$, the oracle returns the
$p$th bit of $E$, defaulting to false beyond its end, and that $E$ has length at most
$B$. For every list $as$ and every register state $p,c,a$ such that $E$ after position
$p$ consists of the encoding of $as$ followed by a suffix, with $c,a\le B$, there is an
$a'\le B$ such that the machine can move from its argument-list scan state with
registers $(p,c,a)$ to the gate scan state with registers
$(p+|\operatorname{encode}(as)|,c,a')$, while all configurations before the endpoint
keep every register at most $B$.
\end{lemma}

\begin{lemma}[Halting during edge-list membership check]
\lean{BoolCircuit.AdjScan.argsEx}
\label{BoolCircuit.AdjScan.argsEx}
\leanok
\uses{BoolCircuit.AdjScan.Gd, BoolCircuit.AdjScan.Target, BoolCircuit.AdjScan.X, BoolCircuit.AdjScan.astep_jeqFst, BoolCircuit.AdjScan.astep_rd, BoolCircuit.AdjScan.drop_add_of_drop, BoolCircuit.AdjScan.drop_succ_of_drop, BoolCircuit.AdjScan.gd_rg, BoolCircuit.AdjScan.getD_of_drop, BoolCircuit.AdjScan.le_of_drop, BoolCircuit.AdjScan.length_encodeList_cons, BoolCircuit.AdjScan.rg, BoolCircuit.AdjScan.rg_zero, BoolCircuit.AdjScan.scan, BoolCircuit.AdjScan.uloop, BoolCircuit.AdjScan.upd_two, BoolCircuit.AdjScan.upd_zero, BoolCircuit.encodeList, BoolCircuit.encodeNat, Complexity.LogProg.AHalt, Complexity.LogProg.AHalt.ret, Complexity.LogProg.AHalt.step, Complexity.LogProg.AReach.halt, Complexity.LogProg.astep}
\difficulty{6}
Let $n$ be a natural number, let $w,E$ be bit strings, and let $o$ be an oracle on bit
strings. Suppose that on every encoding of a natural number $p$, the oracle returns the
$p$th bit of $E$, defaulting to false beyond the end of $E$; suppose also that the input
formed from $n$ and $w$ has edge target $(iv,tv)$. If $E$ has length at most $B$, then
for every register state $p,c,a$ and suffix $r$ such that $E$ after position $p$
consists of the encoded list $as$ followed by $r$, and $c,a\le B$, the edge-scanning
machine started at its argument-list scan state halts with answer true exactly when $iv$
occurs in $as$.
\end{lemma}

\begin{lemma}[Length of a gate encoding]
\lean{BoolCircuit.AdjScan.length_gate_encode}
\label{BoolCircuit.AdjScan.length_gate_encode}
\leanok
\uses{BoolCircuit.DAGGate, BoolCircuit.DAGGate.encode, BoolCircuit.GateKind.encode, BoolCircuit.encodeList, BoolCircuit.encodeNat}
\difficulty{3}
For any Boolean circuit gate $g$, its encoding has length two plus the length of the
list encoding of the natural-number encodings of its argument indices.
\end{lemma}

\begin{lemma}[Bounded gate scan step]
\lean{BoolCircuit.AdjScan.gateSkip}
\label{BoolCircuit.AdjScan.gateSkip}
\leanok
\uses{BoolCircuit.AdjScan.Gd, BoolCircuit.AdjScan.X, BoolCircuit.AdjScan.argsSkip, BoolCircuit.AdjScan.drop_add_of_drop, BoolCircuit.AdjScan.gd_rg, BoolCircuit.AdjScan.le_of_drop, BoolCircuit.AdjScan.length_gate_encode, BoolCircuit.AdjScan.rg, BoolCircuit.AdjScan.rg_one, BoolCircuit.AdjScan.rg_zero, BoolCircuit.AdjScan.scan, BoolCircuit.AdjScan.upd_one, BoolCircuit.AdjScan.upd_zero, BoolCircuit.DAGGate, BoolCircuit.DAGGate.encode, BoolCircuit.GateKind.encode, BoolCircuit.encodeList, BoolCircuit.encodeNat, Complexity.LogProg.AReach, Complexity.LogProg.AReach.step, Complexity.LogProg.astep}
\difficulty{5}
Let $\tau$ be a scan task, let $n$ be a natural number, let $w$ and $E$ be finite bit
strings, and let $o$ assign a Boolean answer to each bit string at the sole oracle
index. Suppose that for every natural number $q$, the oracle answer on the pairing of
the unary string of length $n$ with the binary expansion of $q$ equals the bit of $E$ at
position $q$, taking the answer to be false beyond the end of $E$. Let $B$ be a natural
number with $|E|\le B$, let $g$ be a directed acyclic graph gate, and let $p,c,a$ be
natural numbers such that the suffix of $E$ beginning at position $p$ is the
concatenation of $1$, the encoding of $g$, and a remainder, with $c+1\le B$ and $a\le
B$. Then there is a natural number $a'\le B$ such that the scan for task $\tau$, started
in its gate-reading state with registers $(p,c,a)$ and no answer yet, reaches its next
gate-reading state with registers $(p+1+|\mathrm{encode}(g)|,c+1,a')$ and no answer yet,
with every register at most $B$ in each configuration strictly before arrival.
\end{lemma}

\begin{lemma}[Decoding an encoded gate kind]
\lean{BoolCircuit.AdjScan.decodeKind_encode}
\label{BoolCircuit.AdjScan.decodeKind_encode}
\leanok
\uses{BoolCircuit.AdjScan.decodeKind, BoolCircuit.GateKind, BoolCircuit.GateKind.encode}
\difficulty{4}
For any gate kind $k$, if its encoding is the two-bit string $(b_1,b_2)$, then decoding
those two bits returns $k$.
\end{lemma}

\begin{lemma}[Scan hit on a gate]
\lean{BoolCircuit.AdjScan.hitGate}
\label{BoolCircuit.AdjScan.hitGate}
\leanok
\uses{BoolCircuit.AdjScan.Gd, BoolCircuit.AdjScan.Target, BoolCircuit.AdjScan.X, BoolCircuit.AdjScan.argsEx, BoolCircuit.AdjScan.astep_rd, BoolCircuit.AdjScan.decodeKind_encode, BoolCircuit.AdjScan.gd_rg, BoolCircuit.AdjScan.getD_of_drop, BoolCircuit.AdjScan.hitAns, BoolCircuit.AdjScan.le_of_drop, BoolCircuit.AdjScan.length_gate_encode, BoolCircuit.AdjScan.rg, BoolCircuit.AdjScan.rg_zero, BoolCircuit.AdjScan.scan, BoolCircuit.AdjScan.upd_zero, BoolCircuit.DAGGate, BoolCircuit.DAGGate.encode, BoolCircuit.GateKind.encode, BoolCircuit.encodeList, BoolCircuit.encodeNat, Complexity.LogProg.AHalt, Complexity.LogProg.AHalt.ret, Complexity.LogProg.AHalt.step, Complexity.LogProg.astep}
\difficulty{5}
Let $\tau$ be a task, let $n,iv,tv,B$ be natural numbers, let $w,E$ be finite bit
strings, and let $o$ be an oracle. Suppose that every oracle query on the encoded unary
string of length $n$ paired with the binary expansion of $q$ returns the bit of $E$ at
position $q$, taking out-of-range bits to be $0$. Suppose also that the task target
condition holds for the input encoding the unary string of length $n$ paired with $w$,
that $|E|\le B$, and that the suffix of $E$ at position $p$ begins with $1$ followed by
the encoding of a gate $g$ and a remainder. If $c\le B$ and $a\le B$, then the scan
started in its gate-hit state with registers $(p,c,a)$ and no answer yet halts with
every register bounded by $B$, answering true for a size task, answering whether the
requested gate kind equals the kind of $g$ for a type task, and answering whether $iv$
is among the arguments of $g$ for an edge task.
\end{lemma}

\begin{lemma}[Gate-list scan and halt]
\lean{BoolCircuit.AdjScan.gloop}
\label{BoolCircuit.AdjScan.gloop}
\leanok
\uses{BoolCircuit.AdjScan.Gd, BoolCircuit.AdjScan.Target, BoolCircuit.AdjScan.X, BoolCircuit.AdjScan.astep_hit, BoolCircuit.AdjScan.astep_rd, BoolCircuit.AdjScan.drop_add_of_drop, BoolCircuit.AdjScan.gateSkip, BoolCircuit.AdjScan.gd_rg, BoolCircuit.AdjScan.getD_of_drop, BoolCircuit.AdjScan.hitAns, BoolCircuit.AdjScan.hitGate, BoolCircuit.AdjScan.le_of_drop, BoolCircuit.AdjScan.rg, BoolCircuit.AdjScan.scan, BoolCircuit.DAGGate, BoolCircuit.DAGGate.encode, BoolCircuit.encodeList, Complexity.LogProg.AHalt, Complexity.LogProg.AHalt.ret, Complexity.LogProg.AHalt.step, Complexity.LogProg.AReach.halt}
\difficulty{6}
Let $\tau$ be a task, let $n,iv,tv,B$ be natural numbers, let $w,E$ be finite bit
strings, and let $o$ be an oracle. Suppose that every oracle query on the encoded unary
string of length $n$ paired with the binary expansion of $q$ returns the bit of $E$ at
position $q$, taking out-of-range bits to be $0$. Suppose also that the task target
condition holds for the input encoding the unary string of length $n$ paired with $w$,
the length of $E$ is at most $B$, and $gs$ is a list of gates. For every decomposition
in which the suffix of $E$ at position $p$ is the encoding of the list $gs$ followed by
a remainder, if $c\le tv$, $c+|gs|\le B$, and $a\le B$, then the scan started in its
gate-reading state with registers $(p,c,a)$ and no answer yet halts with every register
bounded by $B$. Its answer is the task-specific answer for the gate of $gs$ at offset
$tv-c$ when $tv-c<|gs|$, and is false otherwise.
\end{lemma}

\begin{lemma}[Initial scan-loop behavior]
\lean{BoolCircuit.AdjScan.nloop}
\label{BoolCircuit.AdjScan.nloop}
\leanok
\uses{BoolCircuit.AdjScan.Gd, BoolCircuit.AdjScan.Target, BoolCircuit.AdjScan.X, BoolCircuit.AdjScan.astep_hit, BoolCircuit.AdjScan.astep_rd, BoolCircuit.AdjScan.bit_rep, BoolCircuit.AdjScan.gd_rg, BoolCircuit.AdjScan.le_of_drop, BoolCircuit.AdjScan.nAns, BoolCircuit.AdjScan.rg, BoolCircuit.AdjScan.rg_one, BoolCircuit.AdjScan.rg_zero, BoolCircuit.AdjScan.scan, BoolCircuit.AdjScan.upd_one, BoolCircuit.AdjScan.upd_zero, Complexity.LogProg.AHalt, Complexity.LogProg.AHalt.ret, Complexity.LogProg.AHalt.step, Complexity.LogProg.astep}
\difficulty{6}
Fix a scan task, natural numbers $n,iv,tv,B$, bit strings $w,E$, an oracle $o$, and a
bit string $rest$. Let $X(n,w)=\langle 1^n,w\rangle$, using the self-delimiting pair
encoding. Suppose the oracle answers each query $X(n,\operatorname{bits}(p))$ with the
$p$th bit of $E$, defaulting to false; the target condition holds, meaning that for the
edge task the two decoded words of $X(n,w)$ are $\operatorname{bits}(iv)$ and
$\operatorname{bits}(tv)$, and for either other task its plain decoded word is
$\operatorname{bits}(tv)$; and $|E|\le B$. Suppose also that $E=1^n 0::rest$, and that
whenever $n\le tv$, the scan machine halts from label $gA$ with registers $(n+1,n,0)$,
respecting the bound $B$ on every register throughout, and returns the bit $ans$. Then
for every $j\le n$ with $j\le tv$, the machine halts from label $nA$ with registers
$(j,j,0)$, respecting the same register bound, and returns $nAns$ if $tv<n$, and
otherwise returns $ans$.
\end{lemma}

\begin{definition}[Task-specific input validity]
\lean{BoolCircuit.AdjScan.Valid}
\label{BoolCircuit.AdjScan.Valid}
\leanok
\uses{BoolCircuit.AdjScan.Task, Complexity.LogProg.ValidPair, Complexity.LogProg.ValidPlain}
For a task $	au$ and a Boolean word $x$, define validity to mean that $x$ is a valid
pair encoding when $	au$ is the edge task, and a valid plain encoding for either other
task.
\end{definition}

\begin{theorem}[Halting of the Circuit Scan]
\lean{BoolCircuit.AdjScan.scan_halts}
\label{BoolCircuit.AdjScan.scan_halts}
\leanok
\uses{BoolCircuit.AdjScan.Gd, BoolCircuit.AdjScan.Target, BoolCircuit.AdjScan.Task, BoolCircuit.AdjScan.Valid, BoolCircuit.AdjScan.X, BoolCircuit.AdjScan.drop_add_of_drop, BoolCircuit.AdjScan.gd_rg, BoolCircuit.AdjScan.gloop, BoolCircuit.AdjScan.hitAns, BoolCircuit.AdjScan.nAns, BoolCircuit.AdjScan.nloop, BoolCircuit.AdjScan.rg, BoolCircuit.AdjScan.rg_init, BoolCircuit.AdjScan.scan, BoolCircuit.DAGCircuit, BoolCircuit.DAGCircuit.encode, BoolCircuit.DAGCircuit.encode_eq, BoolCircuit.DAGCircuit.size, BoolCircuit.DAGCircuit.size_le_length_encode, BoolCircuit.DAGGate.encode, BoolCircuit.encodeList, Complexity.LogProg.AHalt, Complexity.LogProg.AHalt.step, Complexity.LogProg.astep}
\difficulty{4}
Let $C$ be a finite acyclic Boolean circuit with $n$ input vertices, gates listed in
topological order, and an output vertex among its $n+|C|$ vertices, where $|C|$ is the
number of gates. Let $o$ be an oracle whose answer on the binary representation of each
natural number $p$, paired with the unary word of length $n$, is the $p$th bit of the
circuit encoding, with value false when that bit is out of range. Suppose the input word
$w$ encodes a valid instance of the selected task: for the edge task it encodes a pair
of canonical binary words, and for the size or type task it encodes one canonical binary
word. Suppose also that the decoded target index or indices are $i_v$ and $t_v$,
respectively: in the edge task these are the two words’ represented numbers, and
otherwise $t_v$ is the represented number. Then the scan machine for that task, started
with all three registers zero and no answer, halts with all registers at most the
encoding length of $C$. Its answer is the task’s default answer if $t_v<n$; if $t_v\ge
n$ and $t_v-n$ indexes a gate of $C$, its answer is true for the size task, is whether
that gate has the requested kind for the type task, and is whether its argument list
contains $i_v$ for the edge task; if the index is beyond the gate list, its answer is
false.
\end{theorem}
